\subsection{Основные этамы решения}

\begin{enumerate}
    \item Постановка прикладного решения.
    \item Мат. моделирование задачи.
    ...
\end{enumerate}


\section{Часть 1. Элементральная теория погрешонстей}

\subsection{Пункт 1. Источники и теория погрешонстей}

\noindent Источники погрешностей:
\begin{enumerate}
    \item Приближённость мат. моделей \textit{(Любая точная мат. модель содержит погрешность.)}
    \item Погрешность исходных данных.
    \item Приближённость численных методов. \textit{(Численные методы делятся на точные и приближённые, но по факту все явл. приближёнными.)}
    \item Неизбежная потеря точности при представлении чисел в памяти компьютера.
\end{enumerate}

\noindent Источники погрешностей:

\begin{itemize}
    \item Неустранимые \textit{(невозможно устранить методами вычислительной математики)}.
    \item Устранимые \textit{(усовершествовав численные методы, можно приблизить к 0)}.
\end{itemize}

\subsection{Абсолютная и относительная погрешности}

\textbf{Точные значения:} $a, b, c$
\textbf{Приближённые значения:} $a^*, b^*, c^*$

\noindent\textbf{Опр.}: погрешность приближённого значения $a^*$~--- $\epsilon_a = a - a^*$

\noindent\textbf{Опр.}: погрешность абсолютного значения $a$~--- $\epsilon_a = |a - a^*|$

\noindent\textbf{Опр.}: погрешность относительного значения ~--- $\delta a^* = \frac{|a - a^*|}{|a|}$
